\chapter{考察}

\section{結果の解釈}

\subsection{オンデバイス学習による性能改善の特性}

表\ref{tab:results_before}、表\ref{tab:results_after}、および表\ref{tab:improvement}から、
オンデバイス学習により以下の特徴的な変化が観察された：

\begin{itemize}
    \item Accuracy：77.8\% → 85.2\%（+7.3\%pt）
    \item Stability：78.1\% → 85.2\%（+7.1\%pt）
    \item Confidence：84.9\% → 84.4\%（-0.5\%pt）
\end{itemize}

特筆すべきは、\textbf{Confidenceがほぼ不変である一方で、AccuracyとStabilityが顕著に改善している}点である。
これは一見矛盾するように見えるが、実際には本研究のオンデバイス学習戦略（出力層のみを更新）が
適切に機能していることを示す重要な証拠である。

\subsection{判断境界の調整と確信度の分離}

本研究では、ネットワークの出力層（$W_2, \mathbf{b}_2$）のみを更新し、
隠れ層（$W_1, \mathbf{b}_1$）は固定する戦略を採用した。
この設計により、以下の役割分担が実現されている：

\begin{description}
    \item[隠れ層（固定）] 入力特徴量を抽象的な特徴空間に写像する。
    この特徴表現は事前学習により獲得されており、変更されない。
    
    \item[出力層（更新可能）] 特徴空間内でのクラス間の判断境界を定義する。
    オンデバイス学習により、個人の信号特性に合わせて境界を微調整する。
\end{description}

この構造において、Confidenceの変化が小さいことは以下を意味する：

\begin{enumerate}
    \item \textbf{特徴空間の保持}：
    隠れ層が固定されているため、入力信号が写像される特徴空間は変化しない。
    
    \item \textbf{判断境界の平行移動}：
    出力層の重みとバイアスの更新により、判断境界が平行移動または回転するが、
    softmaxの"温度"（確率の鋭さ）は維持される。
    
    \item \textbf{過信の抑制}：
    少量のデータ（約230サンプル）での学習にもかかわらず、
    確率が極端に0.99や1.0に偏ることがない。
    これは過学習・過信を避けた健全な学習状態を示す。
\end{enumerate}

\subsection{低精度ラベルの大幅改善}

表\ref{tab:improvement}から、特にpinchラベルで顕著な改善が見られた
（Accuracy：45.8\% → 70.6\%、+24.8\%pt）。
これは事前学習モデルがpinchジェスチャの認識に苦手意識を持っていたが、
個人の装着条件に適応することで性能が大幅に向上したことを示す。

一方、paperラベルでは若干の性能低下が見られた
（Accuracy：100\% → 96.0\%、-4.0\%pt）。
これは以下の理由が考えられる：

\begin{itemize}
    \item 事前学習時点で既に100\%の精度を達成しており、改善の余地がない
    \item オンデバイス学習により、他のラベル（特にpinch）を識別するために
    判断境界が調整され、paperの領域が若干縮小した
    \item 限られたデータでの学習において、4クラス間のバランスを取る過程で生じた自然なトレードオフ
\end{itemize}

重要なのは、全体平均としては明確な性能向上が見られ、
特に識別困難だったクラスで大幅な改善が得られた点である。

\subsection{実用性能の観点からの評価}

Confidenceの平均値がほぼ不変であることは、実用システムとしては望ましい特性である：

\begin{enumerate}
    \item \textbf{誤検出の抑制}：
    モデルが過剰に自信を持つことがないため、
    間違った予測に高い確率を出すリスクが低い。
    
    \item \textbf{安定した閾値設定}：
    システムがジェスチャ認識を「確定」するための確率閾値を、
    学習前後で変更する必要がない。
    
    \item \textbf{ユーザー体験の一貫性}：
    学習後に突然動作が敏感になったり鈍感になったりすることがなく、
    予測可能な動作を維持できる。
\end{enumerate}

一方、AccuracyとStabilityの向上により：

\begin{itemize}
    \item 正しいジェスチャがより確実に認識される
    \item 予測の時間的な揺らぎが減少し、使用感が向上する
    \item 誤認識による意図しない動作が減少する
\end{itemize}

\subsection{Last-Layer Personalizationの有効性}

本研究の結果は、Transfer Learningの一手法である
Last-Layer Personalization（最終層のみの個人適応）が、
EMG信号のような個人差が大きい生体信号に対して有効であることを示している。

この手法の利点は：

\begin{itemize}
    \item \textbf{学習の効率性}：更新パラメータが100個（全体の27.5\%）に限定され、
    計算コストと学習時間が抑えられる
    
    \item \textbf{汎化性能の維持}：隠れ層の汎用的な特徴抽出能力を保持しつつ、
    個人差に対応できる
    
    \item \textbf{過学習の抑制}：少量データでも安全に学習できる
\end{itemize}

Confidenceが変化しないことは、この設計の妥当性を裏付ける結果と言える。

\section{課題と限界}
本研究で直面した課題や現時点での限界を述べる。

\section{今後の展望}
将来の改善案や応用可能性について記載する。
